\section{Introduction}

Retrieving information from memory strengthens it more than rereading does (the testing effect), and
learning distributed over time is more durable than massed learning (the spacing effect)
\citep{roediger2006,cepeda2006}. Both effects are combined in spaced education, in which short
multiple-choice questions recur periodically over weeks or months. In studies involving physicians and
other health professionals, it improved the acquisition and retention of knowledge
\citep{kerfoot2009,phillips2019}.

In organizations, a similar format is used in \emph{microlearning}, that is, short training units,
often in the form of quizzes. Evidence of its effectiveness, however, comes mainly from studies with
small samples and from grey literature \citep{taylor2022}, and the effects of mandatory employee
training, e.g., in information security, tend to be small or to fade over time
\citep{reinheimer2020}. Training effectiveness is usually assessed through transfer, that is, the
application of knowledge at work \citep{blume2010}. Transfer, however, requires that the knowledge
first be acquired and retained.

This earlier stage can be examined using data from microlearning platforms. The platform records each
response together with its time and outcome, and the same question returns to the same employee
repeatedly. Successive responses are therefore repeated measurements of the same knowledge in the same
person. Similar data are used in learner models applied in intelligent tutoring systems and language
learning applications \citep{corbett1994,settles2016}.

The study was conducted in a bank in which employees answered short questions displayed on their
computers during the working day, and the system collected the responses over three years. Questions
recurred on a schedule independent of response outcome, and the correct answer was displayed after
every response. Such a schedule makes it possible to compare error corrections and forgetting without
the selection introduced by adaptive algorithms, which is rare in data from training platforms.
Laboratory studies further indicate that a selected distractor can become fixed in memory and that
errors made with high confidence can recur despite feedback \citep{butler2008,butler2011}. Whether this
phenomenon occurs in employee training is not known. Three research questions were formulated:
(RQ1) does the percentage of correct responses increase when the same question returns to the same
employee;
(RQ2) how does accuracy depend on the outcome of the previous attempt and on the time elapsed since it;
(RQ3) do repeated errors persist despite feedback, which would point to misconceptions rather than to
guessing.
